% !TeX root = ../papers/nc-p3.tex
\section{薄い二色による文字直線の一意性と印なし同型}
\label{sec:unique-character-frames}
\begin{theorem}\label{thm:unique-character-frames}
前節の重根・単根の両標本を含む開集合で、三通りに絞った四位置の文字直線の
分解は一意である。この開集合の標準正則窓同士の添字保持同型は、32個の置換と
非零係数の積だけで判定できる。二十点配置の安定群の検査を必要としない。
\end{theorem}
\begin{proof}
第一辺の低ブロックの列像の射影直線を使う。厚い四次成分の低ブロックを
$(1,t)^t(1,s)$と置き、高ブロックの行列式を$s$の二次式として判別式を取ると、
射影直線への次数2の射影の分岐式
$B(t)=b_4t^4+b_2t^2+b_0$を得る。奇数冪がないことは真の文字符号の作用からも従う。
二標本で四根は相異なり、定数項と最高次係数の比は順に
\[
 3195075625/15968016,\qquad 9604/184041 .
\]
従って$b_0^2\ne b_4^2$を保つ開集合を取れる。

真の文字の対合はこの直線で$t\mapsto-t$である。残る二元のいずれかを
別の文字分解の符号と仮定し、その行列を$S$とする。前節の四元群で可換で、
真の対合と固定点を共有しないため、$S$は反対角行列である。
新しい文字分解の符号も、薄い鉛筆の二つの文字方向をそれぞれ保つ必要がある。
これらは二つの零点対を同時に交換する対合の固定点として既に一意であり、
右上ブロックでは$I$と交換行列$J$である。高い側の相手の変換を$U$と書けば、
$SIU^{-1}$は$I$の倍数、$SJU^{-1}$は$J$の倍数である。
従って$SJS^{-1}$は$J$の倍数で、その係数は$\pm1$である。
反対角な$S$についてこれは$t\mapsto 1/t$または$-1/t$を与える。
ところが$S$は四次成分を保つので分岐因子も保つ。
そのためには$t^4B(1/t)$が$B(t)$の倍数となり、$b_0^2=b_4^2$が必要で矛盾する。
よって第一辺の一つの文字対しか残らず、系\ref{cor:three-character-frames}で
全四位置への延長も一意である。

文字直線を回復した後、窓同型は各位置で単項行列になる。ブロックを保つ置換は
$ i\mapsto i+u_j+\epsilon_j i_{\rm hi}\,1$の形である。
ここで和は四元群の和、$1$は下位ビットの元である。
テンソルの非零語は総文字零である。高ビットの全ての偶パターンで少なくとも
一つの非零語がある。全て低い場合と全て高い場合は可逆な主ブロック、
連続する同じ組から別組へ移る場合は非零な交差ブロック、他は薄い係数を使う。
四位置の置換後にも総文字零なので、全て低い語から$\sum u_j=0$、
二つずつ高い語から$\epsilon_j$は全位置に共通と分かる。
二つの隣接切断で厚い色の組を保つため$u_0,u_1,u_2$の高ビットが一致し、
和零から$u_3$も同じになる。従って候補数は$2\cdot4\cdot2\cdot2=32$であり、
係数の零を含む標本でも同じ上界である。

各置換後、支持が一致することを確かめる。非零語を行、位置と文字を列とする
整数行列の各行に、その語の四つの文字で1を置く。この行列を$M$、係数比を
$c_\alpha$とすれば、対角基底変換で一致する条件は
\[
 \prod_\alpha c_\alpha^{n_\alpha}=1\qquad(n\in\ker_{\mathbb Z}M^t)
\]
である。Smith標準形で整数核の基底を取り、残る非零対角成分では代数閉体で
根を取れば十分性も従う。全体スカラーは一位置の倍率で吸収できる。
標準正則窓では最上テンソルがスカラーまで決まり、その四切断で関係列が戻る。
定理\ref{thm:canonical-four}により標準四周期を選ぶこの比較が添字保持同型の必要十分判定となる。
\path{computations/verify_unique_character_frames.py}は上の分岐式と32置換を有理数上で検査する。

両標本自身で合格する自己置換は恒等一つである。対角倍率の重み行列の
非零Smith因子は両標本とも1が11個、2が2個である。
従って標準四窓表示の位置ごとのスカラーを除く射影安定群は、
既知の同時文字符号の四元群に一致する。
\end{proof}
